% ============================================================
% 纯答案册·课时09-圆与圆的位置关系 —— 工具/答案抽册器.py 抽册生成件（勿手改；第三档＝纯答案单独成册）
% 源件（只读）：C:/提示词/工作区/M3-第2章量产0913/成卷/练习件/课时09-圆与圆的位置关系/main.tex（md5 9c03a6c71857497a99df15b93a2bd197）
%% 口径：题面不重印；逐题＝题号(印面号同正册)＋[答案]＋[详解]，值/详解照源件逐字
%       （钉值门硬断言：册值≡正册件内值≡值台账值tex，逐字全等）。册序＝正册装配序＝印面号 1..16 升序（同序同号）；键序断言基准＝题面库冻结 manifest canonical 键序。
%% 三档导出：整体 main-true／纯题 main-false（在产）｜本册＝纯答案册（本件）
%% 双档：ansbook-true.tex＝详解印本；ansbook-false.tex＝纯值速查本（详解不印）
% 编译：xelatex <壳> ×2，cwd＝本目录；sty＝qp-m3.sty 本地挂载（md5 7c3930362be8a0a2bdf21bbf8ac16573，锁校验）
% ============================================================
\documentclass[fontset=none]{ctexart}
\usepackage{qp-m3}
% —— 印面逐字符号路由补钉（承源件；− 走 TNR、▱ NSC 子块；禁改 sty 保 md5） ——
\xeCJKDeclareCharClass{Default}{"2212}%
\setCJKfamilyfont{nscblk}[Path=C:/提示词/工作区/字替对照-0909/variantF/fonts/,BoldFont=NSC-w600.ttf]{NSC-w500.ttf}%
\xeCJKDeclareSubCJKBlock{nscblk}{"25B1}%
\newcommand{\pxparallelogram}{{\CJKfamily{nscblk}▱}}%
% —— 断行微调（承母版实测档，三零门承重；\emergencystretch 不另设） ——
\xeCJKsetup{CJKglue={\hskip 0.10em plus 0.08em minus 0.05em}}%
% —— 纯答案册全线括线（长详解可跨栏断，承练习件全线括线口径；灰底盒不可跨栏断） ——
\ansblockgrayfalse
% —— 双档开关（false 档＝纯值速查：答案印、详解不印；件面开关，不动 sty 本体） ——
\newif\ifansbookdetail
\ifdefined\ansbookdetail
  \ifnum\ansbookdetail=0 \ansbookdetailfalse\else\ansbookdetailtrue\fi
\else
  \ansbookdetailtrue
\fi
% —— 页脚件名（承源件章名；件名＝<件型>答案册） ——
\renewcommand{\qpzhangming}{第二章\quad 平面解析几何}%
\renewcommand{\qpjianming}{练习件答案册}%

\begin{document}
\keshi{第9课时\quad 圆与圆的位置关系}
\par\glueguard{1}\addvspace{2pt}\noindent\biaoqian{【纯答案册】}{\fangsong 题面不重印\quad 印面号与正册同号同序}\par\addvspace{2pt}

\begin{multicols}{2}
\emergencystretch=1em

\begin{ansblock}[2章-练-课时09-T1]
% ans:2章-练-课时09-T1
\ansitem{1}{D}
\ifansbookdetail
\ansnote{详解}{两条公切线\(\Longleftrightarrow\)两圆相交．圆心\((0,-m)\)，\((m,0)\)，半径\(\sqrt{2}\)，\(2\sqrt{2}\)，圆心距\(d=\sqrt{2}|m|\)．相交\(\Longleftrightarrow 2\sqrt{2}-\sqrt{2}<\sqrt{2}|m|<2\sqrt{2}+\sqrt{2}\)，即\(1<|m|<3\)．故选：D．}
\fi
\end{ansblock}

\begin{ansblock}[2章-练-课时09-T2]
% ans:2章-练-课时09-T2
\ansitem{2}{C}
\ifansbookdetail
\ansnote{详解}{\(M\cap N=N\Longleftrightarrow N\subseteq M\)，即圆\(N\)内含或内切于圆\(M\)：\(d+r\leqslant 2\)．\(d=\sqrt{2}\)，故\(\sqrt{2}+r\leqslant 2\)，得\(0<r\leqslant 2-\sqrt{2}\)．故选：C．}
\fi
\end{ansblock}

\begin{ansblock}[2章-练-课时09-T3]
% ans:2章-练-课时09-T3
\ansitem{3}{A}
\ifansbookdetail
\ansnote{详解}{\(PA\cdot PB=0\Longleftrightarrow P\)在以\(AB\)为直径的圆\(x^{2}+y^{2}=4\)上（除\(A\)、\(B\)）．存在这样的\(P\Longleftrightarrow\)两圆相交：\(|r-2|<3<r+2\)，得\(r>1\)且\(r<5\)．故选：A．（恰为\(A\)、\(B\)的情形对应相切\(r=1\)或\(r=5\)，已由「不同于\(A\)，\(B\)」排除．）}
\fi
\end{ansblock}

\begin{ansblock}[2章-练-课时09-T4]
% ans:2章-练-课时09-T4
\ansitem{4}{D}
\ifansbookdetail
\ansnote{详解}{满足条件的\(l\)即分别以\(A\)、\(B\)为圆心，1、2为半径的两圆的公切线．圆心距\(|AB|=3=1+2\)，两圆外切，公切线有3条（2外＋1内）．故选：D．}
\fi
\end{ansblock}

\begin{ansblock}[2章-练-课时09-T5]
% ans:2章-练-课时09-T5
\ansitem{5}{C}
\ifansbookdetail
\ansnote{详解}{两方程相减得公共弦\(AB\)：\(x=-2/3\)．\(|AB|=2\sqrt{4-4/9}=8\sqrt{2}/3\)，\(|O_1O_2|=3\)．对角线互相垂直的四边形面积\(=(1/2)\times(8\sqrt{2}/3)\times 3=4\sqrt{2}\)．故选：C．}
\fi
\end{ansblock}

\begin{ansblock}[2章-练-课时09-T6]
% ans:2章-练-课时09-T6
\ansitem{6}{C}
\ifansbookdetail
\ansnote{详解}{设\(P(x_0,y_0)\)．以\(MP\)为直径的圆过切点\(B\)、\(C\)（\(MB\perp PB\)），其方程为\((x-1)(x-x_0)+y(y-y_0)=0\)，与圆\(M\)：\(x^{2}+y^{2}-2x-3=0\)相减得公共弦\(BC\)所在直线：\par\noindent \((x_0-1)x+y_0y-x_0-3=0\)．\(A(2,1)\)在\(BC\)上：\(2(x_0-1)+y_0-x_0-3=0\)，即\(x_0+y_0=5\)．故\(P\)的轨迹为\(x+y-5=0\)．故选：C．}
\fi
\end{ansblock}

\begin{ansblock}[2章-练-课时09-T7]
% ans:2章-练-课时09-T7
\ansitem{7}{ACD}
\ifansbookdetail
\ansnote{详解}{圆心\(M(2,1)\)，\(N(-2,-1)\)，\(d=2\sqrt{5}>2\)，两圆相离，公切线4条．逐项验圆心到直线的距离是否等于1：A：\(y=0\)到\(M\)、\(N\)的距离均为1，是；B：\(3x-4y=0\)到\(M\)的距离\(|6-4|/5=2/5\neq 1\)，不是；C：\(x-2y+\sqrt{5}=0\)到\(M\)、\(N\)的距离均为\(\sqrt{5}/\sqrt{5}=1\)，是；D：同理距离为1，是．故选：ACD．}
\fi
\end{ansblock}

\begin{ansblock}[2章-练-课时09-T8]
% ans:2章-练-课时09-T8
\ansitem{8}{ACD}
\ifansbookdetail
\ansnote{详解}{圆即\((x-3)^{2}+y^{2}=9\)．A：\(r=3\)，正确；B：点\((1,2\sqrt{2})\)到圆心\((3,0)\)的距离\(\sqrt{4+8}=2\sqrt{3}>3\)，在圆外，错误；C：\(d=|3+0+3|/2=3=r\)，相切，正确；D：圆心距\(|3-(-1)|=4\)，半径3、2，\(1<4<5\)，相交，正确．故选：ACD．}
\fi
\end{ansblock}

\begin{ansblock}[2章-练-课时09-T9]
% ans:2章-练-课时09-T9
\ansitem{9}{AD}
\ifansbookdetail
\ansnote{详解}{\(|OC|=\sqrt{4+9}=\sqrt{13}>2+1=3\)，两圆外离，公切线4条，A正确．\par\noindent B：\(x+y=5\)过\(C(2,3)\)且两截距均为5；\(x-y+1=0\)过\(C\)，但两截距为\(-1\)与\(1\)，不相等，故B错误．\par\noindent C：\(C\)在圆\(O\)外（\(|OC|=\sqrt{13}>2\)），过\(C\)的切线应有两条，选项只列一条且表述为唯一，不成立，C错误．\par\noindent D：\(|PQ|\)的最大值为\(|OC|+2+1=\sqrt{13}+3\)，最小值为\(|OC|-2-1=\sqrt{13}-3\)，D正确．故选：AD．}
\fi
\end{ansblock}

\begin{ansblock}[2章-练-课时09-T10]
% ans:2章-练-课时09-T10
\ansitem{10}{AD}
\ifansbookdetail
\ansnote{详解}{四条公切线\(\Longleftrightarrow\)两圆外离．圆心\((0,a)\)，\((a,0)\)，\(d=\sqrt{2}|a|\)，半径3、1．外离\(\Longleftrightarrow\sqrt{2}|a|>4\)，即\(|a|>2\sqrt{2}\)．\(-4\)满足，3满足（\(3>2\sqrt{2}\approx 2.83\)），\(-2\)与\(2\sqrt{2}\)均不满足．故选：AD．}
\fi
\end{ansblock}

\begin{ansblock}[2章-练-课时09-T11]
% ans:2章-练-课时09-T11
\ansitem{11}{1}
\ifansbookdetail
\ansnote{详解}{两方程相减得公共弦所在直线：\(2ay-2=0\)，即\(y=1/a\)，到圆心\(O\)的距离\(d=1/a\)．由半弦长平方\(3=4-1/a^{2}\)，得\(a^{2}=1\)，又\(a>0\)，故\(a=1\)．}
\fi
\end{ansblock}

\begin{ansblock}[2章-练-课时09-T12]
% ans:2章-练-课时09-T12
\ansitem{12}{5}
\ifansbookdetail
\ansnote{详解}{\(\angle APB=90^\circ\Longleftrightarrow P\)在以\(AB\)为直径的圆\(x^{2}+y^{2}=m^{2}\)上．存在这样的\(P\Longleftrightarrow\)两圆有公共点：\(|m-1|\leqslant|OC|=\sqrt{9+7}=4\leqslant m+1\)．由\(4\geqslant m-1\)得\(m\leqslant 5\)（另一侧\(4\leqslant m+1\)给\(m\geqslant 3\)），故\(m\)的最大值为5．}
\fi
\end{ansblock}

\begin{ansblock}[2章-练-课时09-T13]
% ans:2章-练-课时09-T13
\ansitem{13}{相交}
\ifansbookdetail
\ansnote{详解}{\(C_2\)圆心\((2,-m/2)\)在对称轴上：\(2-(\sqrt{3}/2)m+1=0\)，得\(m=2\sqrt{3}\)．\(r_2^{2}=4+3-3=4\)，\(r_2=2\)，圆心\(C_2(2,-\sqrt{3})\)；\(C_1\)圆心\((4,4)\)，\(r_1=5\)．\par\noindent \(d^{2}=4+(4+\sqrt{3})^{2}=23+8\sqrt{3}\approx 36.9\)，而\((r_1-r_2)^{2}=9\)，\((r_1+r_2)^{2}=49\)，\(9<23+8\sqrt{3}<49\)，即\(3<d<7\)，两圆相交．}
\fi
\end{ansblock}

\begin{ansblock}[2章-练-课时09-T14]
% ans:2章-练-课时09-T14
\ansitem{14}{x+2y+1=0}
\ifansbookdetail
\ansnote{详解}{圆即\((x-1)^{2}+(y-1)^{2}=4\)，圆心\((1,1)\)，\(r=2\)．\(S=2\cdot(1/2)r\cdot|MA|=2\sqrt{|MC|^{2}-4}\)，最小\(\Longleftrightarrow|MC|\)最小，即\(C\)到\(l\)的距离\(|1+2+2|/\sqrt{5}=\sqrt{5}\)．\par\noindent 此时\(M\)为垂足：\(CM\)方向\((1,2)\)，直线\(CM\)：\(y-1=2(x-1)\)，与\(l\)联立得\(M(0,-1)\)．以\(CM\)为直径的圆\(x^{2}+y^{2}-x-1=0\)与圆\(C\)相减得\(x+2y+1=0\)，即切点弦\(AB\)的方程．}
\fi
\end{ansblock}

\begin{ansblock}[2章-练-课时09-T15]
% ans:2章-练-课时09-T15
\ansitem{15}{猜想正确，理由见解析}
\ifansbookdetail
\ansnote{详解}{设交点\(A(x_1,y_1)\)，\(B(x_2,y_2)\)，则\(A\)、\(B\)坐标同时满足两圆方程：\(x_i^{2}+y_i^{2}+D_1x_i+E_1y_i+F_1=0\)与\(x_i^{2}+y_i^{2}+D_2x_i+E_2y_i+F_2=0\)（\(i=1,2\)）．\par\noindent 两式相减：\((D_1-D_2)x_i+(E_1-E_2)y_i+(F_1-F_2)=0\)，即\(A\)、\(B\)都在直线\((D_1-D_2)x+(E_1-E_2)y+(F_1-F_2)=0\)上．又\(D_1\neq D_2\)或\(E_1\neq E_2\)（否则两圆方程相同），该直线系数不全为零，而过两个不同交点的直线有且仅有一条，故它就是公共弦\(AB\)所在直线的方程，猜想正确．}
\fi
\end{ansblock}

\begin{ansblock}[2章-练-课时09-T16]
% ans:2章-练-课时09-T16
\ansitem{16}{(x−4)²+(y−3)²=25}
\ifansbookdetail
\ansnote{详解}{设圆心\(D(a,b)\)．切点\(B\)在两圆连心线上（连心线过切点且与公切线垂直）：\(k_{CB}=k_{DB}\)，即\(b/a=6/8=3/4\)，\(b=3a/4\)．又\(|DA|=|DB|\)：\((a-1)^{2}+(b+1)^{2}=(a-8)^{2}+(b-6)^{2}\)，展开得\(14a+14b=98\)，即\(a+b=7\)．联立\(b=3a/4\)得\(a=4\)，\(b=3\)，\(r=|DB|=\sqrt{16+9}=5\)．\par\noindent 故所求圆的方程为\((x-4)^{2}+(y-3)^{2}=25\)．（验：\(|DA|=\sqrt{9+16}=5\)；\(|DC|=10-5=5\)，两圆内切于\(B\)．）}
\fi
\end{ansblock}

% ---- 尾块（\tailfill[30mm] 定高参——钉档 §三.3：无参在平衡栏呈「内容尾随框」，贴栏底须定高参；实测无参致末页平衡 Overfull\vbox 12.99pt，定高 30mm 三零） ----
\tailfill[30mm]

\end{multicols}
\end{document}
